\textbf{Momento angular, producto vectorial y sistema de ecuaciones.} Primero se calculan las tres filas de $\boldsymbol H=I\boldsymbol\omega$. Las entradas $xz$ y $zx$ son negativas con el convenio del enunciado.\[\fitmath{\boldsymbol H=\left[\begin{matrix}927.0\\-1040.0\\2475.6\end{matrix}\right]\,\mathrm{kg\,m^2/s}}\]\[\fitmath{H_{1}=(8500)(0.12)+(-620)(0.15)\approx927\,\mathrm{kg\,m^2/s}}\]\[\fitmath{H_{2}=(13000)(-0.08)\approx-1040\,\mathrm{kg\,m^2/s}}\]\[\fitmath{H_{3}=(-620)(0.12)+(17000)(0.15)\approx2475.6\,\mathrm{kg\,m^2/s}}\]\[\fitmath{c_{1}=(-0.08)(2475.6)-(0.15)(-1040)\approx-42.048\,\mathrm{N\,m}}\]\[\fitmath{c_{2}=(0.15)(927)-(0.12)(2475.6)\approx-158.022\,\mathrm{N\,m}}\]\[\fitmath{c_{3}=(0.12)(-1040)-(-0.08)(927)\approx-50.64\,\mathrm{N\,m}}\]El producto giroscópico se desarrolla como $(qH_z-rH_y,\ rH_x-pH_z,\ pH_y-qH_x)$; su valor es:\[\fitmath{\boldsymbol\omega\times\boldsymbol H=\left[\begin{matrix}-42.048\\-158.022\\-50.64\end{matrix}\right]\,\mathrm{N\,m}}\]Se resta de los momentos aplicados, obteniendo el segundo miembro $\boldsymbol b=\boldsymbol M-\boldsymbol\omega\times\boldsymbol H$.\[\fitmath{\boldsymbol b=\left[\begin{matrix}992.048\\-341.978\\390.64\end{matrix}\right]\,\mathrm{N\,m}}\]La ecuación de $y$ queda desacoplada: $\dot q=b_y/I_{yy}$. Para $x,z$, eliminar la otra aceleración en el sistema $I_{xx}\dot p-I_{xz}\dot r=b_x$, $-I_{xz}\dot p+I_{zz}\dot r=b_z$ da:\[\fitmath{\dot p=\frac{I_{zz}b_x+I_{xz}b_z}{I_{xx}I_{zz}-I_{xz}^2},\qquad\dot r=\frac{I_{xz}b_x+I_{xx}b_z}{I_{xx}I_{zz}-I_{xz}^2}}\]\[\fitmath{(\dot p,\dot q,\dot r)=\left[\begin{matrix}0.118703407542279 & -0.026306 & 0.0273080066280125\end{matrix}\right]\,\mathrm{s^{-2}}}\]Para comprobar, multiplicar el vector de aceleraciones por $I$ y sumar el término giroscópico recupera $\boldsymbol M$. El trabajo giroscópico se anula por perpendicularidad; la tasa de energía es $\boldsymbol\omega\cdot\boldsymbol M=205\,\mathrm W$.
